\documentclass[12pt]{article}
\usepackage[utf8]{inputenc}
\usepackage[spanish]{babel}
\usepackage{amsmath}
\usepackage{amsfonts}
\usepackage{amssymb}
\usepackage{geometry}
\geometry{letterpaper, margin=2.5cm}

\begin{document}

\begin{center}
    \Large \textbf{Evaluación de Examen 1: Unidad I} \\
    \large Física para Ciencias de la Salud (005-1713) \\
    \vspace{0.5cm}
    \textbf{Estudiante:} Ybramaris Maribao \quad \textbf{C.I.:} 36.150.392
    \vspace{0.3cm} \\
    \large \textbf{Calificación Final: 5/10 puntos}
\end{center}

\vspace{0.5cm}

\section*{Análisis de Resultados}

\subsection*{1. Ángulo plano (Conversión a radianes)}
\textbf{Procedimiento y resultado:} Plantea correctamente la conversión multiplicando por el factor $\frac{\pi}{180^\circ}$. Evidencia de manera explícita la simplificación de la fracción dividiendo el numerador y denominador entre 15, obteniendo de forma directa el resultado exacto de $\frac{5\pi}{12}$ Rad. \\
\textbf{Veredicto:} \textbf{Correcto} (2/2 pts).

\subsection*{2. Sistemas de coordenadas (Longitud del hueso)}
\textbf{Procedimiento y resultado:} Excelente resolución. Dibuja el gráfico identificando de forma clara los catetos. Aplica el Teorema de Pitágoras llevando un excelente control de las unidades ($l^2 = 36\text{ cm}^2 + 64\text{ cm}^2 = 100\text{ cm}^2$) hasta despejar correctamente y obtener $10$ cm. \\
\textbf{Veredicto:} \textbf{Correcto} (2/2 pts).

\subsection*{3. Vectores (Fuerza resultante)}
\textbf{Procedimiento y resultado:} Presenta un error conceptual respecto a la física del problema. El enunciado solicita la fuerza resultante, lo cual es la sumatoria de fuerzas ($\vec{F}_1 + \vec{F}_2$). La estudiante plantea y calcula una diferencia de vectores ($\Delta\vec{F} = \vec{F}_2 - \vec{F}_1$). Aunque la operación matemática de su resta es correcta, no responde a lo requerido. \\
\textbf{Veredicto:} \textbf{Parcialmente correcto} (Manejo correcto del álgebra, pero error conceptual en el planteamiento de la resultante) (1/2 pts).

\subsection*{4. Potencias de diez (Notación científica y prefijo)}
\textbf{Procedimiento y resultado:} Presenta errores matemáticos en el desplazamiento de la coma. Expresa la notación científica como $1,25 \times 10^{-4}$ m (el exponente correcto es $-5$). Además, asocia de manera equivocada la cifra "1,25" con el prefijo "micrómetros", lo que no coincide con el valor original (el correcto era $12,5$ micrómetros). \\
\textbf{Veredicto:} \textbf{Incorrecto} (0/2 pts).

\subsection*{5. Sistemas de unidades (Conversión de densidad)}
\textbf{Procedimiento y resultado:} La estudiante dejó esta pregunta completamente en blanco. \\
\textbf{Veredicto:} \textbf{Incorrecto} (0/2 pts).

\vspace{0.5cm}
\section*{Comentarios Generales y Sugerencias}
La estudiante demuestra poseer habilidades sólidas en aritmética básica y geometría plana.
\begin{itemize}
    \item Se recomienda repasar con atención los conceptos teóricos: la "fuerza resultante" es una suma (Pregunta 3) y se debe tener mayor cuidado al contar los espacios decimales en la notación científica (Pregunta 4).
    \item Sugerencia importante: Siempre administrar el tiempo del examen para intentar algún planteamiento en todas las preguntas y evitar dejar ejercicios en blanco (Pregunta 5).
\end{itemize}

\end{document}